\section{Paying for capacity: direct settlement, and a service built over it}
\label{sec:paying-for-capacity}

Buying capacity on \Flux is closer to buying a train ticket than to opening a utility account. A
tenant computes what a deployment costs, sends that many FLUX to a known address in one transaction,
and every node independently reads the chain and agrees that the deployment is funded until a
particular block height. Nobody issues a credential, nobody keeps a ledger of what the tenant owes,
and when the paid-for height arrives the entitlement simply stops. This section specifies that
mechanism first, because it is the protocol's settlement path and the one autonomous software should
use. It then describes something that is not part of the protocol at all: a centralized payment
service, operated by InFlux Technologies on top of \Flux
\srcintent{design intake}, which holds a person's keys, signs for them, and pays the
chain on their behalf in exchange for a card payment. The chain does not know that service exists,
and this section proves it: from consensus's point of view its payments are ordinary payments from
ordinary keys, and nothing in \fluxd or \FluxOS is aware that a custodian is involved
(Property~\ref{prop:chainindifferent}). Last, and separately, it specifies --- explicitly as unimplemented --- the deposited-balance layer
such a service would need in order to become a credit account. The three are kept apart throughout,
because they belong to different systems, carry different maturities, and fail in different ways.
The claim the section defends is narrow and falsifiable: \textbf{\Flux's settlement path is trustless
and ships today in two independent forms; a third party has built a convenience layer over it for
people who want to pay by card, and the protocol neither knows about that layer nor depends on it.}

\paragraph{Symbols introduced here.}
Not in \texttt{notation.tex}, used only in this section: $e_\app$ the paid-through height of
application $\app$; $L^{\ast}(\hgt)$ the effective default lifetime; $\iota(\hgt)$ the active
rate-table row; $\varpi$ the unused fraction of a previous term; $v^{\mathrm{obs}}$ an observed
on-chain payment. For the third-party service of Section~\ref{sec:paying:facilitator}: $\mathcal{F}$
the operator's facilitator service; $K_u$ the pair of private keys it holds in trust for user $u$ and
never discloses to $u$; $\mathcal{H}$ its set
of operator hot wallets; $T_{\mathrm{rec}}$ its record retention time; $\Delta_{\mathrm{adv}}$ its
early-renewal advance window in blocks; $\mathsf{A}_{\mathrm{orb}}, \mathsf{A}_{\mathrm{ren}},
\mathsf{A}_{\mathrm{paid}}$ its three hardcoded allowlists. For the unimplemented credits layer of
Section~\ref{sec:paying:creditspec}: $b_u$ a tenant's deposited balance in $\sat$;
$\mathrm{st}(\app,\hgt)$ the lifecycle state; $G_{\mathrm{g}}, G_{\mathrm{s}}$ the grace and
suspension windows in blocks; $\Delta_{\mathrm{pre}}$ the pre-authorisation lead in blocks;
$\price^{\mathrm{ren}}$ and $\price^{\mathrm{res}}$ the renewal and resumption charges;
$\Delta^{\mathrm{sub}}$ the suspension subsidy. The second group is candidate notation for an object
that does not exist and must not be absorbed into \texttt{notation.tex} unless it is built.

% =====================================================================
\subsection{Part A --- Direct settlement}
\label{sec:paying:direct}

Everything in this part is \shipped and is written in the present tense, and everything in it is a
property of the network itself --- what \fluxd and \FluxOS do, with no service standing in between.
It describes two deployed instances of the same idea: the application-hosting path in \FluxOS, and
the \CumulusVPN entitlement path. Both share one property --- the payment is the entitlement --- and
neither maintains any forward-looking claim against the payer.

\subsubsection{The price function and the rate schedule}

\begin{assumption}[Settlement environment]
\label{asm:settle}
As stated in the system model of Section~2: (i) every node holds the same confirmed chain prefix up
to the reorg bound (\texttt{MAX\_REORG\_LENGTH} $=40$ outside the PoN stabilisation window);
(ii) every node runs \FluxOS with the same compiled rate table, subject to on-chain
\texttt{chainparams} price messages of version \texttt{p} posted by the Foundation multisig, which
take effect from their occurrence height \srccode{flux/\allowbreak{}ZelBack/\allowbreak{}src/\allowbreak{}services/\allowbreak{}utils/\allowbreak{}chainUtilities.js}{9--52};
(iii) application payments are recognised by output address plus an \texttt{OP\_RETURN} message
carrying the temporary-message hash \srccode{flux/\allowbreak{}ZelBack/\allowbreak{}src/\allowbreak{}services/\allowbreak{}explorerService.js}{309--345}.
Assumption (ii) is a trust assumption on the Foundation and is carried as such in Section~7
(the price table is the base; consensus fixes only the split).
\end{assumption}

\begin{definition}[Height-scheduled rate table]
\label{def:ratetable}
The rate table is a finite sequence of segments indexed by activation height,
$\{(h_j, \pricevec^{(j)})\}_{j=0}^{5}$ with $h_0 = -1$ (genesis) and $h_1 < \dots < h_5$ as given in
Table~\ref{tab:pricetable}. Each row is
$\pricevec^{(j)} = \bigl(p_{\mathrm{cpu}}, p_{\mathrm{ram}}, p_{\mathrm{hdd}}, p_{\min},
p_{\mathrm{port}}, p_{\mathrm{scope}}, p_{\mathrm{sip}}\bigr)^{(j)} \in \mathbb{Q}_{>0}^{7}$,
denominated in FLUX. The active row at height $\hgt$ is $\pricevec(\hgt) \defeq
\pricevec^{(\iota(\hgt))}$ with $\iota(\hgt) \defeq \max\{\, j : h_j < \hgt \,\}$ (strict: the
code selects rows with \texttt{height < height}
\srccode{flux/\allowbreak{}ZelBack/\allowbreak{}src/\allowbreak{}services/\allowbreak{}utils/\allowbreak{}appUtilities.js}{32}).
\end{definition}

\begin{table}[H]
\centering\small
\caption{The deployment rate schedule: six segments, genesis plus five revisions. Values as
returned by the live endpoint and as compiled into \FluxOS.
\srcmeasured{api.runonflux.io/\allowbreak{}apps/\allowbreak{}deploymentinformation}{2026-09-03};
\srccode{flux/\allowbreak{}ZelBack/\allowbreak{}config/\allowbreak{}default.js}{133--209}. \shipped}
\label{tab:pricetable}
\begin{tabular}{@{}lrrrrrrr@{}}
\toprule
segment & height $h_j$ & $p_{\mathrm{cpu}}$ & $p_{\mathrm{ram}}$ & $p_{\mathrm{hdd}}$ &
$p_{\min}$ & $p_{\mathrm{port}}$ & $p_{\mathrm{scope}}$ \\
\midrule
genesis     & $-1$              & 3    & 1    & 0.5   & 1    & 2   & 6   \\
revision 1  & \num{983000}      & 0.3  & 0.1  & 0.05  & 0.1  & 2   & 6   \\
revision 2  & \num{1004000}     & 0.06 & 0.02 & 0.01  & 0.01 & 2   & 6   \\
revision 3  & \num{1288000}     & 0.15 & 0.05 & 0.02  & 0.01 & 2   & 6   \\
revision 4  & \num{1594832}     & 0.15 & 0.05 & 0.02  & 0.01 & 1.5 & 6   \\
revision 5  & \num{1597156}     & 0.03 & 0.01 & 0.004 & 0.01 & 0.4 & 0.8 \\
\bottomrule
\end{tabular}
\end{table}

Across the six segments the CPU rate falls from 3 to 0.03, a factor of \num{100}; the RAM rate falls
from 1 to 0.01 and the disk rate from 0.5 to 0.004. The static-IP surcharge $p_{\mathrm{sip}}$ is 3
in segments 0--4 and 0.4 in segment 5. A seventh segment is staged at height \num{2950000} with
values identical to segment 5, so that \specv{9} registrations price against a defined segment
\srcbranch{flux}{feat/appspec-v9}{ZelBack/config/default.js:196--209}; it is not a revision and is not counted here
(conflict C27). A parallel USD-denominated table and a \texttt{/\allowbreak{}apps/\allowbreak{}getappspecsusdprice} read
endpoint also exist in configuration \srccode{flux/\allowbreak{}ZelBack/\allowbreak{}config/\allowbreak{}default.js}{211--224}; which
denomination governs is a Foundation policy lever, discussed as such in Section~7, and this paper
makes no claim about it.

\begin{definition}[Deployment price]
\label{def:price}
Let $\app$ be an application with component set $C(\app)$ and per-component resources
$r_c = (r^{\mathrm{cpu}}_c, r^{\mathrm{ram}}_c, r^{\mathrm{hdd}}_c) \in \mathbb{Q}_{\ge 0} \times
\mathbb{N} \times \mathbb{N}$, in vCPU, \unit{\mega\byte} and \unit{\giga\byte} respectively, and let
$\resvec(\app) \defeq \sum_{c \in C(\app)} r_c$. Let $\ninst \in \{k_{\min}(\hgt), \dots, 100\}$ be
the instance count, $\dur \in \{d_{\min}(\hgt), \dots, \num{1056000}\}$ the duration in blocks, and
$\hgt \in \mathbb{N}$ the height at which the registration is priced. Write
$\lceil x \rceil_2$ for rounding up to two decimals and $\mathcal{E}(\app)$ for the set of
enterprise ports the specification requests. The \emph{base}, in FLUX per default lifetime, is
\begin{multline*}
B(\app,\hgt) \;\defeq\;
10\,p_{\mathrm{cpu}}\,\resvec^{\mathrm{cpu}}
\;+\; \tfrac{1}{100}\,p_{\mathrm{ram}}\,\resvec^{\mathrm{ram}}
\;+\; p_{\mathrm{hdd}}\,\resvec^{\mathrm{hdd}} \\
{}\;+\; p_{\mathrm{scope}}\,\indic{\text{targeted or enterprise}}
\;+\; p_{\mathrm{sip}}\,\indic{\text{static IP}}
\;+\; p_{\mathrm{port}}\,\lvert\mathcal{E}(\app)\rvert ,
\end{multline*}
all rates evaluated at $\pricevec(\hgt)$
\srccode{flux/\allowbreak{}ZelBack/\allowbreak{}src/\allowbreak{}services/\allowbreak{}utils/\allowbreak{}appUtilities.js}{89--122}. The base is quoted per
three instances: write $\beta(\app,\hgt) \defeq \lceil B(\app,\hgt)/3 \rceil_2$ for the per-instance
base. The charge per default lifetime is
\[
A(\app,\ninst,\hgt) \;\defeq\;
\begin{cases}
\beta, & \ninst = 1 \text{ or } \hgt < \num{1890000},\\[2pt]
\beta + 0.5\,(\ninst - 1), & \beta < 0.5 \text{ and } \ninst > 3,\\[2pt]
\bigl\lceil \beta\,(\ninst - 1)\bigr\rceil_2 + \beta, & \text{otherwise,}
\end{cases}
\]
\srccode{flux/\allowbreak{}ZelBack/\allowbreak{}src/\allowbreak{}services/\allowbreak{}utils/\allowbreak{}appUtilities.js}{123--134}
\srccode{flux/\allowbreak{}ZelBack/\allowbreak{}config/\allowbreak{}default.js}{306}; the USD-minimum branch that follows it, gated by
\texttt{apply\allowbreak{}Minimum\allowbreak{}Price\allowbreak{}On3Instances} \srccode{flux/\allowbreak{}ZelBack/\allowbreak{}src/\allowbreak{}services/\allowbreak{}utils/\allowbreak{}appUtilities.js}{136--138},
is inert on the FLUX path, since no row of the FLUX rate table carries \texttt{minUSDPrice}
\srccode{flux/\allowbreak{}ZelBack/\allowbreak{}src/\allowbreak{}services/\allowbreak{}utils/\allowbreak{}chainUtilities.js}{22--36}. The price of a fresh
registration is
\[
\price(\app,\ninst,\dur,\hgt) \;=\;
\max\Bigl\{\,p_{\min}(\hgt),\;
\Bigl\lceil\; \frac{\dur}{L^{\ast}(\hgt)} \cdot A(\app,\ninst,\hgt) \;\Bigr\rceil_2 \Bigr\} ,
\qquad
L^{\ast}(\hgt) \defeq \begin{cases}
\blockslasting, & \hgt < \hpon,\\[2pt]
4\,\blockslasting, & \hgt \ge \hpon,
\end{cases}
\]
with $\blockslasting = \num{22000}$ blocks \srccode{flux/\allowbreak{}ZelBack/\allowbreak{}config/\allowbreak{}default.js}{285} and
$\hpon = \num{2020000}$ \srccode{flux/\allowbreak{}ZelBack/\allowbreak{}config/\allowbreak{}default.js}{293}
\srccode{flux/\allowbreak{}ZelBack/\allowbreak{}src/\allowbreak{}services/\allowbreak{}appMessaging/\allowbreak{}messageVerifier.js}{736--745}. An \emph{update} is
priced at $\max\{p_{\min},\, 0.9\,(\price - \varpi\,\price_{\mathrm{prev}})\}$, where
$\varpi = (\dur_{\mathrm{prev}} - (\hgt - \hgt_{\mathrm{prev}}))/\dur_{\mathrm{prev}}$ is the unused
fraction of the previous term
\srccode{flux/\allowbreak{}ZelBack/\allowbreak{}src/\allowbreak{}services/\allowbreak{}appMessaging/\allowbreak{}messageVerifier.js}{797--828}.
\end{definition}

\noindent\textbf{Bounds and constants.}
$\blockslasting = \num{22000}$ blocks, which at the pre-\PoN{} \qty{2}{\minute} spacing is
\qty{30.6}{\day} \srccode{flux/\allowbreak{}ZelBack/\allowbreak{}config/\allowbreak{}default.js}{285};
$k_{\min} = 3$ in general and $1$ for \specv{8} applications from height \num{2176519}
\srccode{flux/\allowbreak{}ZelBack/\allowbreak{}config/\allowbreak{}default.js}{257--259}; $k_{\max} = 100$
\srccode{flux/\allowbreak{}ZelBack/\allowbreak{}config/\allowbreak{}default.js}{260}; the post-PoN duration ceiling is \num{1056000} blocks
($\approx$ \qty{367}{\day} at \qty{30}{\second}) \srccode{flux/\allowbreak{}ZelBack/\allowbreak{}config/\allowbreak{}default.js}{292} and
the floor $d_{\min}(\hgt)$ is height-gated by the validator: \num{5000} blocks, then \num{100} from height
\num{1630040}, then \num{1} from height \num{1964447}
\srccode{flux/\allowbreak{}ZelBack/\allowbreak{}src/\allowbreak{}services/\allowbreak{}appRequirements/\allowbreak{}appValidator.js}{852--862},
while the live endpoint still reports \num{5000}
\srcmeasured{api.runonflux.io/\allowbreak{}apps/\allowbreak{}deploymentinformation}{2026-09-03} (conflict C33).
The payment sink is the transparent address \texttt{t3Nryf\allowbreak{}AQLGe\allowbreak{}Fs9j\allowbreak{}Eoeqsxm\allowbreak{}BN2QLRa\allowbreak{}RKFLUX} from height
\num{1670000}, preceded by \texttt{t3a\allowbreak{}GJvdtd8NR6Grnqn\allowbreak{}Ru\allowbreak{}VEz\allowbreak{}H6Mbr\allowbreak{}Xu\allowbreak{}JFLUX} and, at any height, the
legacy \texttt{t1LUs6quf7TB2z\allowbreak{}VZmexq\allowbreak{}PQdnqmr\allowbreak{}FMGZGj\allowbreak{}V6}
\srccode{flux/\allowbreak{}ZelBack/\allowbreak{}config/\allowbreak{}default.js}{241--245}. All application revenue therefore lands at a
single transparent address by policy, which Section~2 records as a centralization surface and
Section~7 replaces with a consensus-enforced split. Measured, \num{439} of \num{1195} deployed
applications (\qty{36.7}{\percent}) sit at the three-instance floor and the maximum observed instance
count is exactly $k_{\max}$ \srcmeasured{api.runonflux.io/\allowbreak{}apps/\allowbreak{}globalappsspecifications}{2026-09-03}.
The minimum is re-applied \emph{after} the duration
scaling \srccode{flux/\allowbreak{}ZelBack/\allowbreak{}src/\allowbreak{}services/\allowbreak{}appMessaging/\allowbreak{}messageVerifier.js}{740--745},
so no registration settles below $p_{\min}$, which is why the floor sits outside the duration factor
in Definition~\ref{def:price}.

\begin{definition}[Expiry and entitlement]
\label{def:expiry}
Let $\hgt_0$ be the height of the block carrying the confirmed payment that promotes $\app$'s
temporary message to a permanent one, and $\dur$ its requested duration. The paid-through height is
\[
e_\app \;\defeq\;
\begin{cases}
\hgt_0 + \dur, & \hgt_0 \ge \hpon,\\[2pt]
\hpon + 4\,(\hgt_0 + \dur - \hpon), & \hgt_0 < \hpon \le \hgt_0 + \dur,
\end{cases}
\]
the second case being the post-fork recomputation applied to applications straddling PoN activation
\srccode{flux/\allowbreak{}ZelBack/\allowbreak{}src/\allowbreak{}services/\allowbreak{}appMessaging/\allowbreak{}messageVerifier.js}{711--724}. Entitlement is the
predicate $\mathrm{Ent}(\app, \hgt) \defeq \indic{\hgt \le e_\app}$, evaluated independently by every
node.
\end{definition}

\subsubsection{The property that makes direct settlement the default}

\begin{property}[No standing obligation]
\label{prop:nostanding}
\shipped. Under direct settlement, for every application $\app$ and every height $\hgt$, the value
$\mathrm{Ent}(\app,\hgt)$ is a function of the confirmed chain prefix at or before $\hgt$ alone.
Consequently: (i) the network holds no forward-looking claim against the tenant --- there exists no
height $\hgt' > \hgt$ at which the tenant owes anything; (ii) for all $\hgt' > \hgt$, the value of
$\mathrm{Ent}(\app,\hgt')$ is already determined at $\hgt$ up to the arrival of a further payment,
and is non-increasing in $\hgt'$ absent one; (iii) tenant and operator cannot hold divergent views of
the entitlement except by holding divergent views of the chain.
\end{property}

\begin{proof}
By Definition~\ref{def:expiry}, $e_\app$ is determined by $(\hgt_0, \dur)$, both of which are fields
of, or the height of, a block at or before $\hgt$; by Definition~\ref{def:price}, the promotion test
$v^{\mathrm{obs}} \ge \price(\app,\ninst,\dur,\hgt_0)$ is evaluated at promotion time against the
rate row $\iota(\hgt_0)$, and $\iota$ is monotone in height, so no later block changes it. Hence
$\mathrm{Ent}(\app,\cdot)$ factors through the prefix, giving (i) and (ii). For (iii): under
Assumption~\ref{asm:settle}(i) two honest nodes agree on the prefix, and $e_\app$ is computed by the
same integer arithmetic on both, so any disagreement is a disagreement about the chain and surfaces
as such.
\end{proof}

Property~\ref{prop:nostanding} is the reason this paper recommends direct settlement to autonomous
software. An agent needs no account, no stored credential and no relationship that outlives the
deployment; there is no balance to track, no charge to authorise, and no reconciliation step that can
fail while the agent is not watching. It is also the reason this path carries no grace, suspension or
eviction machinery: expiry is arithmetic, not a decision, so there is nothing for a policy to decide
and no state a policy could get wrong.

\begin{corollary}
\label{cor:noeviction}
No grace or suspension state is required on this path. Any such state would be keyed on a balance
or an outstanding obligation, and by Property~\ref{prop:nostanding}(i) neither exists.
\end{corollary}

\begin{construction}[\textsc{settle-direct} --- registration, payment, promotion, \shipped]
\label{alg:settle}
\emph{Inputs.} A formatted specification $\app$; a signature over it by the owner's ZelID; a duration
$\dur$; a funded wallet.\quad
\emph{State.} Per node: a temporary-message store (TTL \qty{3600}{\second}), a permanent-message
store, and the chain.
\begin{enumerate}\setlength{\itemsep}{1pt}
\item The tenant reads $\pricevec(\hgt)$ from \texttt{/\allowbreak{}apps/\allowbreak{}deploymentinformation} and computes
      $\price$ by Definition~\ref{def:price}.
\item The tenant \texttt{POST}s the signed message to any node's \texttt{/\allowbreak{}apps/\allowbreak{}appregister}; the node
      verifies the signature, the timestamp window ($\le$ \qty{1}{\hour} old, $\le$ \qty{5}{\minute}
      ahead), the schema for the declared specification version, and its own peer counts
      ($\ge 8$ outgoing, $\ge 4$ incoming), then broadcasts the message and returns only its hash
      \srccode{flux/\allowbreak{}ZelBack/\allowbreak{}src/\allowbreak{}services/\allowbreak{}appDatabase/\allowbreak{}registryManager.js}{1714--1888}.
\item The tenant broadcasts one transaction paying $v^{\mathrm{obs}} \ge \price$ to the sink address,
      carrying that hash in an \texttt{OP\_RETURN}.
\item Each node's explorer loop recognises the payment by output address, sums \texttt{valueSat}
      across matching outputs, decodes the hash, and calls the promotion routine
      \srccode{flux/\allowbreak{}ZelBack/\allowbreak{}src/\allowbreak{}services/\allowbreak{}explorerService.js}{309--345,444}.
\item Promotion stores the permanent message, recomputes $e_\app$ (Definition~\ref{def:expiry}), and
      admits the application to global state iff $v^{\mathrm{obs}} \ge \price$
      \srccode{flux/\allowbreak{}ZelBack/\allowbreak{}src/\allowbreak{}services/\allowbreak{}appMessaging/\allowbreak{}messageVerifier.js}{619--897}.
\end{enumerate}
\emph{Invariant.} After promotion, global state contains $\app$ with paid-through height $e_\app$
computed from the chain alone; before promotion it contains nothing about $\app$ that survives the
temporary-message TTL.\quad
\emph{Failure modes.} (a) The hash is not resolvable from peers within three attempts spaced
\qty{60}{\second}: the registration never promotes and the tenant sees a client-side error at step~2.
(b) $v^{\mathrm{obs}} < \price$: see Property~\ref{prop:silentunderpay} --- the payment is spent and
the tenant is told nothing.
\end{construction}

\subsubsection{\CumulusVPN: the worked example}

The second deployed instance of direct settlement is not an application-hosting charge at all, and it
is the more important of the two for this paper's thesis. \CumulusVPN is a WireGuard exit-gateway
product whose gateways run as ordinary \Flux applications. Its entitlement mechanism is this: the
client generates a WireGuard keypair on the device, derives a payment code
$\mathsf{code} = \mathsf{base58}\bigl(\mathrm{sha256}(pk)[0{:}20]\bigr)$ from its own public key, and
sends FLUX to a fixed address with the \texttt{OP\_RETURN} memo $\texttt{CVPN1:}\,\mathsf{code}$.
Gateway and client compute the code byte-identically
\srccode{cumulusvpn/\allowbreak{}gateway/\allowbreak{}internal/\allowbreak{}entitle/\allowbreak{}entitle.go}{99--108}
\srccode{cumulusvpn/\allowbreak{}clients/\allowbreak{}core-ts/\allowbreak{}src/\allowbreak{}paymentCode.ts}{17--23}. Every gateway then scans the chain
by itself and derives the map $\mathsf{paidUntil}[\mathsf{code}]$ from confirmed, correctly-memoed
payments, with no gateway-to-gateway coordination of any kind
\srccode{cumulusvpn/\allowbreak{}gateway/\allowbreak{}internal/\allowbreak{}entitle/\allowbreak{}entitle.go}{1--5,218--287}.

\textbf{There is no account, no token and no server-issued credential.} The WireGuard public key
\emph{is} the credential: the gateway checks the key it already sees on the tunnel against the map it
derived from the chain \srccode{cumulusvpn/\allowbreak{}gateway/\allowbreak{}internal/\allowbreak{}api/\allowbreak{}api.go}{246,374}. Authentication runs
in the other direction only --- the control API signs every response body with an ed25519 key derived
from the gateway's own WireGuard private key, so a client that learned the public key through
discovery can verify the responder without a TLS certificate
\srccode{cumulusvpn/\allowbreak{}gateway/\allowbreak{}internal/\allowbreak{}api/\allowbreak{}api.go}{563--585}. One payment covers every gateway,
including both hops of a multi-hop path, because entitlement is keyed by the client's key on all of
them \srcdoc{cumulusvpn/docs/11-multihop.md:47--58}. This is a deployed instance of ``payment is the
account'', which the roadmap lists as a future item \srcroadmap{Q4 2026}; it is not a plan, it is
running (conflict C19).

\begin{construction}[\textsc{derive-entitlement} --- the gateway chain scanner, \shipped]
\label{alg:cvpn}
\emph{Inputs.} A transaction source exposing \texttt{BlockCount} and \texttt{AddressTxs} (the host
node's daemon proxy, falling back to a public explorer); the payment address $A$; the price $P$ in
FLUX per \qty{30}{\day} period; an optional snapshot at \texttt{/\allowbreak{}data/\allowbreak{}entitle.\allowbreak{}state}.\quad
\emph{State.} $\mathsf{paidUntil} : \mathsf{code} \to \text{time}$, plus a scan cursor.
\begin{enumerate}\setlength{\itemsep}{1pt}
\item Backfill from the stored cursor, or from height $0$ if there is none.
\item Every \qty{15}{\second}, read \texttt{BlockCount}; on a new height, fetch the new transactions
      paying $A$.
\item Accept a transaction iff it carries exactly one \texttt{CVPN1:} memo and
      $30\,(v + 10^{-9}) \ge P$, i.e.\ at least one day's worth
      \srccode{cumulusvpn/\allowbreak{}gateway/\allowbreak{}internal/\allowbreak{}entitle/\allowbreak{}entitle.go}{270--287,296--323}.
\item Grant $\lfloor 30 v / P \rfloor$ days, stacking onto any existing
      $\mathsf{paidUntil}[\mathsf{code}]$, capped at $\text{now} + \qty{720}{\day}$
      \srccode{cumulusvpn/\allowbreak{}gateway/\allowbreak{}internal/\allowbreak{}entitle/\allowbreak{}entitle.go}{34--36,218--268}.
\item Checkpoint immediately on a real grant, otherwise at most once per \qty{60}{\second}; treat the
      snapshot as pure cache and discard it on any version, address or price mismatch
      \srccode{cumulusvpn/\allowbreak{}gateway/\allowbreak{}internal/\allowbreak{}entitle/\allowbreak{}snapshot.go}{14,29--31}.
\item Every \qty{15}{\second}, re-derive each enrolled peer's tier and retune its rate limiter in
      place, so a payment takes effect on in-flight flows with no reconnect
      \srccode{cumulusvpn/\allowbreak{}gateway/\allowbreak{}internal/\allowbreak{}limiter/\allowbreak{}limiter.go}{105--123}.
\end{enumerate}
\emph{Invariant.} $\mathsf{paidUntil}$ is a pure function of the confirmed transaction history of $A$;
two gateways with the same view of the chain and the same $(A,P)$ derive identical maps without
exchanging a message. Changing $P$ invalidates every cached snapshot by construction (step~5).\quad
\emph{Failure modes.} If the route to chain data is absent, the engine distinguishes ``route absent''
from ``no payments'' and retains the previous map rather than demoting every paying peer
\srcdoc{cumulusvpn/docs/03-gateway.md:135--138}. If \texttt{/data} is unwritable, the snapshot and the
peer table are lost on restart; a statically issued configuration then cannot re-enroll
\srccode{cumulusvpn/\allowbreak{}gateway/\allowbreak{}internal/\allowbreak{}wg/\allowbreak{}peercache.go}{17--24,47--51}. A migration of the gateway
instance to another node changes the server key and address, which invalidates an already-issued
configuration --- unavoidably so \srcdoc{cumulusvpn/docs/03-gateway.md:95--97}.
\end{construction}

\paragraph{The caveat that travels with it.}
Payment and usage are \textbf{linkable}. The memo is a one-way hash of the public key, but the
gateway a user enrolls with learns the real public key, recomputes the hash, and can then find the
paying address on a public chain. The repository states this itself as a known, unmitigated gap, with
payment-key indirection and blind-signature vouchers named as the intended fixes
\srcdoc{cumulusvpn/docs/04-payments.md:81--87}. \CumulusVPN is therefore simultaneously the best
existing demonstration of direct settlement and a concrete, shipped instance of the problem
Section~21 takes up: the network must be convinced that a workload is funded while learning neither
payer nor workload. Section~21 is written against this system, not against a hypothetical one.

\subsubsection{The defect on this path, stated as a requirement}

\begin{property}[Silent underpayment]
\label{prop:silentunderpay}
\shipped. If a registration's confirmed payment satisfies $v^{\mathrm{obs}} < \price$, the promotion
routine emits a log warning and drops the registration. No error is surfaced to the payer, no retry
is scheduled, and no refund path exists in the code
\srccode{flux/\allowbreak{}ZelBack/\allowbreak{}src/\allowbreak{}services/\allowbreak{}appMessaging/\allowbreak{}messageVerifier.js}{735--763}; the update path
behaves identically \srccode{flux/\allowbreak{}ZelBack/\allowbreak{}src/\allowbreak{}services/\allowbreak{}appMessaging/\allowbreak{}messageVerifier.js}{797--828}.
The funds are spent and the deployment does not exist.
\end{property}

\begin{proof}[Proof sketch]
By inspection of the promotion routine: the price comparison is the last gate before
\texttt{insert\allowbreak{}App\allowbreak{}Specifications}, and its false branch is a \texttt{log.warn} with no message emitted
on the P2P channel and no HTTP response outstanding to write to --- the tenant's only synchronous
interaction ended at step~2 of Construction~\ref{alg:settle}. From the payer's side the observable
state is indistinguishable from ``the message never propagated''.
\end{proof}

For a human tenant this is a support ticket. For an autonomous agent with nobody to escalate to it is
an unrecoverable silent failure, and it sits on exactly the path this paper recommends agents use.
That asymmetry, not the size of the loss, is what makes it the most important correctness gap in
Part~A.

\begin{remark}[Requirement R16.1 --- signed negative acknowledgement]
\label{req:nack}
Any node that declines to promote a registration or update for insufficient payment shall emit a
signed message
$\mathsf{nack} = (\mathsf{hash}, \hgt, \mathsf{txid}, \price^{\mathrm{req}}, v^{\mathrm{obs}},
\mathsf{reason})$
over the existing signed-envelope P2P channel, retrievable by hash from any node's read API. Three
properties are required and one is explicitly \emph{not}: the message must be verifiable against the
deterministic node list under the same rule as every other broadcast; it must carry both the computed
price and the observed payment, so the tenant can locate the discrepancy without a second round trip;
and it must be emitted on the same condition that triggers the drop, so that absence of a
\textsf{nack} is informative. Quorum is \emph{not} required: because $\price$ is a deterministic
public function of the chain (Definition~\ref{def:price}), a single valid \textsf{nack} from any
confirmed node is actionable evidence that the tenant can re-derive itself. The construction is
cheap; its absence is what makes Invariant~\ref{inv:nosilent} below fail on Part~A today.
\end{remark}

% =====================================================================
\subsection{Part B --- A service built on \Flux, and the credits layer that lives outside it}
\label{sec:paying:credits}

This part contains two kinds of object, a reader must not merge them, and \textbf{neither of them is
part of \Flux}. Section~\ref{sec:paying:facilitator} is \shipped and is written in the present tense
because it runs: a centralized payment service, operated by InFlux Technologies
\srcintent{design intake}, through which a person signs in with an email address,
pays by card, and has a deployment registered, funded and renewed for them. It is not a credit
account. It holds no user balance, offers no withdrawal, and keeps no record for longer than six
days; what it does hold is both of the user's private keys, which the user never sees
\srcintent{design intake}, and that is the subject of the subsection. Nearly every
fact stated there carries code provenance and the one that does not says so; and every one of them is
a fact about the operator's service rather than about the network, because the chain cannot observe
the service and does not depend on it (Property~\ref{prop:chainindifferent}). Section~\ref{sec:paying:creditspec} is \planned and is
written in the conditional: the deposited balance and lifecycle state machine a credits layer would
require, retained here as a specification because it is the object Section~\ref{sec:migration}'s dependency list waits
on, and marked throughout as unimplemented. Nothing in the second subsection describes the first.

\begin{remark}[Requirement R16.2 --- inherited from Section~7]
\label{req:pnr}
Section~7 exports one hard requirement to this section: after PNR activation at $\hpa$, every
application-hosting charge, \emph{whatever the payment channel}, must settle as exactly one PNR
transaction on L1, or the consensus-enforced split is void. A credit balance may therefore not become
an off-chain ledger of hosting charges; each charge it authorises must appear on-chain as one
settlement. This constrains the construction below at every point and is the reason
Invariant~\ref{inv:conservation} is stated as a conservation law rather than as bookkeeping. A second
requirement runs the other way: any balance forfeited on eviction must not be routed into the fee
pool $\pool$, or evicting a tenant becomes directly profitable for the fleet that votes (Section~23).
The operator's service already satisfies the first requirement by construction --- it settles each
card charge as exactly one transaction to the sink address (Construction~\ref{alg:prepay}, step~4)
--- which is worth recording, because it means a convenience layer of this shape does not obstruct
PNR. That is a statement about one operator's implementation, not a guarantee the protocol extracts
from any service built over it; the requirement binds the settlement shape, and any service that
violates it simply fails to produce a funded deployment.
\end{remark}

\subsubsection{What runs on top: a third-party prepay-and-forward payment facilitator}
\label{sec:paying:facilitator}

\begin{assumption}[Environment of a third-party payment service]
\label{asm:custodial}
In addition to Assumption~\ref{asm:settle}: (i) the facilitator is a single administrative party
operating server-side infrastructure outside the \FluxNode set, so Section~2's adversary model does
not cover it --- a node honest under that model stays honest while the facilitator misbehaves, and
the converse; (ii) the facilitator authenticates users by email PIN or by Google/Apple OAuth through
a hosted identity provider \srccode{console-api/\allowbreak{}src/\allowbreak{}services/\allowbreak{}auth/\allowbreak{}pin.ts}{4--7}
\srccode{fluxos-frontend/\allowbreak{}src/\allowbreak{}utils/\allowbreak{}firebase.js}{53--142}, and that authentication is orthogonal to
every key the network verifies; (iii) no consensus rule, on-chain record or contract binds the
facilitator's behaviour --- all of it is server-side code and operational discipline; (iv) the
network is not a party to the arrangement and takes no view of it: the accept/reject decision of
Construction~\ref{alg:settle} is invariant under the substitution of the facilitator for the tenant,
so no node can tell that a service was involved (Property~\ref{prop:chainindifferent}, proved below).
Clause (iii) is a trust assumption of exactly the kind Section~2 inventories, and it is carried in
Section~23 as a centralization surface \emph{of this service}; clause (iv) is why it is not a
centralization surface of the network.
\end{assumption}

\begin{definition}[Third-party custodial payment facilitator]
\label{def:facilitator}
Write $\mathrm{addr}(\cdot)$ for the address derived from a public key and $\mathrm{owner}(\app)$ for
the address carried in $\app$'s \texttt{owner} field. A \emph{facilitator} $\mathcal{F}$ is a
commercial operator --- here InFlux Technologies \srcintent{design intake} ---
running server-side infrastructure outside the \FluxNode set and holding, for each user $u$ it
serves, two private keys generated server-side at account creation and never disclosed to $u$
\srccode{console-api/\allowbreak{}src/\allowbreak{}services/\allowbreak{}flux/\allowbreak{}keypair.ts}{20--27}
\srcintent{design intake}:
\begin{itemize}\setlength{\itemsep}{1pt}
\item an \emph{identity} key $sk^{\mathrm{id}}_u$ over secp256k1, whose P2PKH address
$\mathrm{addr}(pk^{\mathrm{id}}_u)$ --- the user's ZelID --- is written into the \texttt{owner} field
of every specification $\mathcal{F}$ registers for $u$, and which signs the off-chain register
message;
\item a \emph{payment} key $sk^{\mathrm{pay}}_u$ on the \Flux transparent network, used to build and
broadcast on-chain transactions.
\end{itemize}
Write $K_u = (sk^{\mathrm{id}}_u, sk^{\mathrm{pay}}_u)$. $K_u$ is held \emph{in trust}: it is
generated by the operator, it stays with the operator, and $u$ is never shown either member of it
\srcintent{design intake}. Both members are stored
AES-256-GCM-encrypted at rest and WIF-encoded, and both are decrypted in the serving process
whenever a signature or a transaction is required
\srccode{console-api/\allowbreak{}src/\allowbreak{}db/\allowbreak{}models/\allowbreak{}FluxIdentity.ts}{13--28}; at most one active identity exists per
account, enforced by a unique index \srccode{console-api/\allowbreak{}src/\allowbreak{}db/\allowbreak{}models/\allowbreak{}FluxIdentity.ts}{40--43}.
$\mathcal{F}$ additionally holds a set $\mathcal{H} = \{h_{\mathrm{fiat}}, h_{\mathrm{free}},
h_{\mathrm{ren}}\}$ of three operator hot wallets, from which every registration and renewal it
performs is funded. Its durable per-payment record is
\[
r \;=\; \bigl(H(m_\app),\; v^{\mathrm{sat}},\; v^{\mathrm{usd}},\; t\bigr)
\;\in\; \{0,1\}^{256} \times \mathbb{Z}_{\ge 0} \times \mathbb{Q}_{\ge 0} \times \mathbb{N},
\]
keyed by \emph{application-message hash} rather than by user, carrying the satoshi amount settled on
chain and the USD amount charged to the card, and subject to a time-to-live index
$T_{\mathrm{rec}} = \qty{518400}{\second} = \qty{6}{\day}$ on every collection
\srccode{appsmonitor/\allowbreak{}config/\allowbreak{}default.js}{33}
\srccode{appsmonitor/\allowbreak{}src/\allowbreak{}services/\allowbreak{}appsFreeUpdatesMonitor.js}{67--71}. \textbf{There is no further
state.} In particular the record carries no user-held quantity, and the balance $b_u$ and lifecycle
label of Definition~\ref{def:creditstate} below have no referent here
(Property~\ref{prop:nobalance}).
\end{definition}

\begin{definition}[Prepay-and-forward settlement]
\label{def:prepayforward}
Let $\mathrm{bal} : \mathcal{H} \to \mathbb{Z}_{\ge 0}$ give the operator wallets' balances in
$\sat$. For a registration or renewal of $\app$ for user $u$ at height $\hgt$, the settlement action
of $\mathcal{F}$ is: charge $u$'s card instrument, held at the payment processor and not at \Flux,
for the corresponding USD amount; then broadcast one transaction from some $h \in \mathcal{H}$,
\[
\mathrm{Pay}_{\mathcal{F}}(u,\app,\hgt):\quad
\mathrm{bal}(h) \;\longmapsto\; \mathrm{bal}(h) - \price(\app,\ninst_\app,\dur_\app,\hgt);
\]
then acknowledge the charge to the processor as fulfilled. The debit falls on the \emph{operator's}
wallet. No user-held quantity is decremented, because none exists. This is the precise sense in
which the service is a \emph{prepay-and-forward payment facilitator} rather than a credit
account: value moves card~$\to$~operator and operator~$\to$~chain, never
card~$\to$~user~balance~$\to$~chain. Only the second leg is visible on chain, and it is visible as an
ordinary payment from an ordinary key.
\end{definition}

\begin{definition}[Discretionary free hosting]
\label{def:freehosting}
$\mathcal{F}$ also funds some applications from $\mathcal{H}$ with no card charge at all. Eligibility
is a predicate $\mathrm{Free}$ on the set of deployed applications, evaluated by the operator's own
daemon and by nothing else,
\[
\mathrm{Free}(\app) \;\defeq\;
\indic{\app \in \mathsf{A}_{\mathrm{orb}}} \;\vee\;
\indic{\mathrm{owner}(\app) \in \mathsf{A}_{\mathrm{ren}}} \;\vee\;
\indic{\mathrm{owner}(\app) \in \mathsf{A}_{\mathrm{paid}}} \;\vee\;
\mathrm{Cap}(\app),
\]
where $\mathsf{A}_{\mathrm{orb}}$ is a hardcoded list of nine grandfathered application names
\srccode{appsmonitor/\allowbreak{}src/\allowbreak{}services/\allowbreak{}appsMonitor.js}{36--46}; $\mathsf{A}_{\mathrm{ren}}$ a hardcoded
owner allowlist whose applications are renewed from $h_{\mathrm{ren}}$ once the adjusted expiry
falls within four days \srccode{appsmonitor/\allowbreak{}config/\allowbreak{}default.js}{60--72}
\srccode{appsmonitor/\allowbreak{}src/\allowbreak{}services/\allowbreak{}appsMonitor.js}{308,437--445}; $\mathsf{A}_{\mathrm{paid}}$ a
hardcoded two-address allowlist whose applications are subsidised in full and indefinitely
\srccode{appsmonitor/\allowbreak{}config/\allowbreak{}default.js}{92--94}
\srccode{appsmonitor/\allowbreak{}src/\allowbreak{}services/\allowbreak{}appsPaidForOwnersMonitor.js}{78--90}; and $\mathrm{Cap}(\app)$
holds iff $\app$ runs the single image \texttt{runonflux/\allowbreak{}orbit:latest} in one component at one
instance with $r^{\mathrm{cpu}} \le 0.5$ vCPU, $r^{\mathrm{ram}} \le \qty{1000}{\mega\byte}$ and
$r^{\mathrm{hdd}} \le \qty{5}{\giga\byte}$, and its owner runs no other such application
\srccode{appsmonitor/\allowbreak{}src/\allowbreak{}services/\allowbreak{}appsMonitor.js}{137--148,159--180}. A separate rule funds a
zero-price update --- one whose price delta is zero and whose requested expiry extension is at most
11 blocks --- at a flat \num{0.02}\flux from $h_{\mathrm{free}}$, after a \qty{2}{\minute} settle
delay \srccode{appsmonitor/\allowbreak{}src/\allowbreak{}services/\allowbreak{}appsFreeUpdatesMonitor.js}{133--137,157,174}
\srccode{appsmonitor/\allowbreak{}src/\allowbreak{}services/\allowbreak{}fluxService.js}{364--367}; and a further rule pays one month for
applications matching a marketplace listing
\srccode{appsmonitor/\allowbreak{}src/\allowbreak{}services/\allowbreak{}appsMarketplaceFreeFirstMonth.js}{342--344}. Every set and
constant above is operator-side configuration: none is on chain, none is user-settable, and no
appeal path exists.
\end{definition}

\begin{construction}[\textsc{prepay-forward} --- custodial registration, payment and renewal, \shipped]
\label{alg:prepay}
\emph{Inputs.} A user $u$ authenticated to $\mathcal{F}$ under Assumption~\ref{asm:custodial}(ii); a
specification $\app$; a card instrument, either a one-off checkout session or a stored
subscription.\quad
\emph{State.} $K_u$ and $\mathcal{H}$ as in Definition~\ref{def:facilitator}; the per-payment records
under $T_{\mathrm{rec}}$; the chain.
\begin{enumerate}\setlength{\itemsep}{1pt}
\item $\mathcal{F}$ decrypts $sk^{\mathrm{id}}_u$ in process and signs $m_\app$, whose \texttt{owner}
      field is $\mathrm{addr}(pk^{\mathrm{id}}_u)$
      \srccode{console-api/\allowbreak{}src/\allowbreak{}db/\allowbreak{}models/\allowbreak{}FluxIdentity.ts}{13--28}. In the legacy stack the
      equivalent step is a remote call that returns both a ZelID address and a signature over a
      network-issued login phrase, which \FluxOS's own \texttt{verifylogin} then accepts
      \srccode{fluxos-frontend/\allowbreak{}src/\allowbreak{}@core/\allowbreak{}components/\allowbreak{}Login.vue}{1039--1104}.
\item The signed message is submitted to a node exactly as in step~2 of
      Construction~\ref{alg:settle}; the node returns $H(m_\app)$ and nothing else.
\item The card is charged in USD.
\item A polling loop enumerates message hashes for which a card charge exists and an on-chain
      payment does not, computes $\price$ by Definition~\ref{def:price}, broadcasts \emph{one}
      transaction from $h_{\mathrm{fiat}}$ to the deployment sink carrying $H(m_\app)$ in
      \texttt{OP\_RETURN}, and acknowledges the charge
      \srccode{appsmonitor/\allowbreak{}src/\allowbreak{}services/\allowbreak{}appsPaidWithFiatMonitor.js}{93--190}
      \srccode{appsmonitor/\allowbreak{}src/\allowbreak{}services/\allowbreak{}fluxService.js}{345--362}.
\item Promotion proceeds exactly as in steps~4--5 of Construction~\ref{alg:settle}. No node
      distinguishes this payment from any other.
\item \emph{Renewal.} Every \qty{3600}{\second} the facilitator reads each active subscription's
      application from global state, computes the \PoN-adjusted remaining lifetime, and when
      $e_\app - \hgt < \Delta_{\mathrm{adv}} = \num{8640}$ blocks ($\approx$ \qty{3}{\day} at
      \qty{30}{\second}) triggers a card charge; the pending-payment loop then repeats steps~4--5
      every \qty{60}{\second}
      \srccode{appsmonitor/\allowbreak{}src/\allowbreak{}services/\allowbreak{}appsSubscriptionRenewalMonitor.js}{112--368,379--460,566--568}.
      The renewed term defaults to \num{88000} blocks and is capped at \num{1056000}
      \srccode{appsmonitor/\allowbreak{}src/\allowbreak{}services/\allowbreak{}appsSubscriptionRenewalMonitor.js}{18,213}.
\end{enumerate}
\emph{Invariant.} At every height, $\mathrm{Ent}(\app,\hgt)$ (Definition~\ref{def:expiry}) is the
same function of the confirmed chain prefix as on Part~A. The facilitator changes \emph{who signs}
and \emph{who pays}; it does not change how the network decides, and it cannot, because every step
above is a step of Construction~\ref{alg:settle} executed by a different party
(Property~\ref{prop:chainindifferent}).\quad
\emph{Failure modes.} (a) The card charge fails: a loop reads the processor's failure list every
\qty{21600}{\second} and, with a per-subscription cooldown of \qty{24}{\hour}, emails the
application's contact address; that is the entire remedy, and no code path revives an application
once expired \srccode{appsmonitor/\allowbreak{}src/\allowbreak{}services/\allowbreak{}appsSubscriptionRenewalMonitor.js}{19,463--575}.
(b) A renewal loop stops or is stopped: it fails closed to \emph{unrenewed}, and the application
expires by Definition~\ref{def:expiry} exactly as an unfunded direct-settlement one would.
(c) $t > T_{\mathrm{rec}}$: the on-chain payment persists and is public, but the operator-side
association between a card charge and an application ceases to exist.
(d) Underpayment: Property~\ref{prop:silentunderpay} applies unchanged, with the loss falling on
$h_{\mathrm{fiat}}$ rather than on the user \srcinferred.
\end{construction}

\begin{property}[Custodial signing and spending authority]
\label{prop:custody}
\shipped. For every user $u$ served by $\mathcal{F}$, the facilitator can (i) produce a valid
signature under $pk^{\mathrm{id}}_u$ over any message --- including any \texttt{appregister},
\texttt{appupdate} or expire message naming $u$'s address in the \texttt{owner} field --- and
(ii) spend any output controlled by $pk^{\mathrm{pay}}_u$; in both cases without a counter-signature
from $u$ and without $u$ being online. No threshold, policy engine, hardware element or second
factor stands between the serving process and either key. The correct description of $u$'s position
is therefore not that $u$ delegates authority to $\mathcal{F}$, but that \textbf{$\mathcal{F}$ is
the key holder and $u$ authenticates to it.}
\end{property}

\begin{proof}[Proof sketch]
For the current stack, by direct inspection: both private keys live in one record, encrypted under a
service-held master key with associated data the service itself supplies, and are decrypted in the
same worker that signs and broadcasts
\srccode{console-api/\allowbreak{}src/\allowbreak{}db/\allowbreak{}models/\allowbreak{}FluxIdentity.ts}{13--28}. Encryption at rest binds use of the key
to possession of the master key, not to any act of $u$; a key-provider abstraction with a
self-hosted KMS backend exists and hardens that binding without changing it \inprogress
\srccode{console-api/\allowbreak{}src/\allowbreak{}lib/\allowbreak{}crypto/\allowbreak{}keys.ts}{63--286}. For the legacy stack the conclusion follows
from the observed protocol rather than from a read of the signer, which is not in the corpus: the
login path sends a network-issued login phrase to a remote service and receives back both a ZelID
address and a signature over that phrase which \texttt{verifylogin} accepts
\srccode{fluxos-frontend/\allowbreak{}src/\allowbreak{}@core/\allowbreak{}components/\allowbreak{}Login.vue}{1039--1104}; producing that signature
requires the corresponding private key, so the service holds it \srcinferred.
\end{proof}

\begin{corollary}[There is nothing to delegate from, so Section~17 does not reach this path]
\label{cor:nodelegation}
\shipped. On this path the delegation object of Section~17 --- a principal $\principal$ granting an
agent $\agent$ a scope $\scopeset$ under a budget $\budget$ --- has no instantiation, because
$\principal$ never holds the key from which authority would be delegated. The correct reading is not
that Section~17's mechanism is present but unenforced here; it is that the mechanism has nowhere to
attach, and no certificate construction can bound an authority its putative grantor does not hold.
Section~19 identifies the same architecture behind its operator-signing path and observes that a
bearer token there is bounded by nothing; Property~\ref{prop:custody} sharpens that: the token does
not confer authority, it authenticates the user to the party that already holds it. Whatever bounds
that path is therefore operator-side policy, not a protocol mechanism, and this paper specifies none
for it. Section~19's local-signing path --- \SSP and Zelcore, where the signature is produced on the
user's own device --- is the path Section~17 is written for, and every property proved there
continues to hold there.
\end{corollary}

\begin{property}[No balance, no withdrawal, no durable record]
\label{prop:nobalance}
\shipped. In the deployed facilitator: (i) no field of the data model denotes a user-held quantity
of FLUX or of fiat, the record being keyed per application-message hash rather than per user
(Definition~\ref{def:facilitator}); (ii) every collection carries a TTL index of
$T_{\mathrm{rec}} = \qty{6}{\day}$, so no payment record older than six days exists
\srccode{appsmonitor/\allowbreak{}config/\allowbreak{}default.js}{33}; (iii) a targeted search of the three repositories that
implement the path found no code path returning FLUX or fiat to a user, hence no withdrawal
\srccode{appsmonitor, fluxos-frontend, console-api}{targeted search, no match}; (iv) the newer
stack's billing endpoint aggregates an on-chain wallet balance with two read-only legacy sources and
marks unconfirmed deposits \texttt{shadow:\allowbreak{}true, credited:\allowbreak{}false} so that they are never summed into a
spendable balance, and exposes no write path at all \inprogress
\srccode{console-api/\allowbreak{}src/\allowbreak{}routes/\allowbreak{}billing.ts}{44--61,63--165}. Consequently $b_u$ of
Definition~\ref{def:creditstate} has no deployed referent, and neither has any state in
$\{\textsf{funded},\textsf{grace},\textsf{suspended},\textsf{evicted}\}$.
\end{property}

\unsupported{A repository named \texttt{withdrawalservice} exists, with a generic claim-intake schema
(\texttt{coin}, \texttt{network}, \texttt{address}, optional \texttt{fluxid}) that reads as a
multi-product claim path rather than as this layer's ledger
\srccode{withdrawalservice/\allowbreak{}src/\allowbreak{}controllers/\allowbreak{}claimController.js}{1--58}. Whether it is wired to the
facilitator at all was not determined. Property~\ref{prop:nobalance}(iii) is the result of a negative
search across the three implementing repositories and must be re-checked against that service before
publication.}

\begin{property}[Payer and owner are independent]
\label{prop:payersep}
\shipped. The promotion test of Construction~\ref{alg:settle} is a function of the paying
transaction's \emph{outputs} only: recognition is by output address and by the temporary-message
hash decoded from \texttt{OP\_RETURN}, and the amount is the sum of \texttt{valueSat} over matching
outputs \srccode{flux/\allowbreak{}ZelBack/\allowbreak{}src/\allowbreak{}services/\allowbreak{}explorerService.js}{309--345}. No input, and no property
of the spending key, is examined. Hence for any application $\app$, the party that pays
$\price(\app,\ninst,\dur,\hgt)$ need not be the party whose key signs $m_\app$ and appears in the
\texttt{owner} field.
\end{property}

\begin{proof}[Proof sketch]
The register message is authenticated by the owner's signature at step~2 of
Construction~\ref{alg:settle} and carries no payer field; the payment is authenticated by nothing
beyond being confirmed on chain, and is bound to the message only through the hash in its
\texttt{OP\_RETURN}. The two authentications are therefore over disjoint objects, and the promotion
routine composes them without a cross-check. The converse capability --- a third party altering an
owner's expiry --- exists only through the separate hardcoded \texttt{usersToExtend} allowlist
\srccode{flux/\allowbreak{}ZelBack/\allowbreak{}config/\allowbreak{}default.js}{229}, which is a distinct mechanism and is not needed here.
\end{proof}

Property~\ref{prop:payersep} is why the custody of Property~\ref{prop:custody} is not forced by the
protocol. Fronting the FLUX for a tenant --- the economically useful half of what $\mathcal{F}$ does
--- requires only the ability to broadcast one transaction naming a message hash. Holding the
tenant's identity key --- the half that carries the trust --- is what card checkout and email login
currently require, not what payment requires. The two halves are separable in the deployed protocol
today; nothing in the corpus separates them. The same fact, read the other way, gives the property
that fixes this subsection's place in the paper.

\begin{property}[The chain is indifferent to the service]
\label{prop:chainindifferent}
\shipped. Fix an application $\app$, a duration $\dur$ and a height $\hgt_0$. Consider two executions
of Construction~\ref{alg:settle}: one in which $u$ signs $m_\app$ on their own device and broadcasts
the paying transaction, and one in which $\mathcal{F}$ signs under $K_u$ and broadcasts from
$h_{\mathrm{fiat}} \in \mathcal{H}$ (Construction~\ref{alg:prepay}). If both signatures verify
against $\mathrm{owner}(\app)$ and both payments satisfy $v^{\mathrm{obs}} \ge \price$, then every
node's observable behaviour is identical in the two executions: the same message is stored, the same
$e_\app$ is computed, and $\mathrm{Ent}(\app,\cdot)$ is the same function. Consequently, across the
registration, promotion, pricing and entitlement paths --- which are the whole of what a node does
with an application payment --- no consensus rule, no \FluxOS validation step and no piece of node
state distinguishes a registration made through a payment service from one made by a tenant alone.
The service is invisible to the protocol, and the protocol's guarantees neither improve nor degrade
in its presence.
\end{property}

\begin{proof}[Proof sketch]
Three inputs, each already established. (a) Promotion is a function of the paying transaction's
outputs only --- recognition by output address, amount by summing \texttt{valueSat} over matching
outputs, binding by the hash decoded from \texttt{OP\_RETURN} --- and examines no input and no
property of the spending key (Property~\ref{prop:payersep}
\srccode{flux/\allowbreak{}ZelBack/\allowbreak{}src/\allowbreak{}services/\allowbreak{}explorerService.js}{309--345}). So the identity of the broadcaster
is not an input to any node's decision. (b) The register message is admitted on verification of a
signature against the address in the \texttt{owner} field
\srccode{flux/\allowbreak{}ZelBack/\allowbreak{}src/\allowbreak{}services/\allowbreak{}appDatabase/\allowbreak{}registryManager.js}{1714--1888}; signature
verification is a predicate on $(\text{message}, \text{signature}, \text{address})$ and takes no
argument describing who holds the corresponding private key, so it cannot separate the two
executions. (c) $\price$ and $e_\app$ are functions of $(\app, \ninst, \dur, \hgt_0)$ alone
(Definitions~\ref{def:price} and~\ref{def:expiry}), and none of those four is a property of the payer
or the signer. The two executions therefore agree on every input to every node-side decision, hence
on every output. The argument is a statement about what the deployed code reads, not about what it
intends to permit; it would be falsified by any node-side branch on payer identity, and none exists
in the cited paths.
\end{proof}

This is the property that fixes what the rest of the subsection is about. The service of
Definition~\ref{def:facilitator} is not a \Flux capability with a trust caveat attached; it is a
company's product sitting on top of a settlement path that would work exactly the same way if the
company shut it down tomorrow. Everything enumerated below is therefore a fact about that product.
None of it is a fact about \Flux, and a reader assessing the network's trust model should discount
all of it.

\paragraph{The trust surface of the operator's service, enumerated.}
Each item is a property of the operator's service only. None of it is a property of \Flux, and none
of it applies to Part~A. Items~1--5 are \shipped on code provenance; item~6 rests on a stated
behaviour and says so.

\begin{enumerate}\setlength{\itemsep}{2pt}
\item \textbf{The operator can sign and spend as the user.} Property~\ref{prop:custody}. There is no
counter-signature, no bound on what may be signed, and no user-side revocation --- revoking would
require holding a key the user does not have.

\item \textbf{The operator decides unilaterally who is hosted for free.}
Definition~\ref{def:freehosting}: three hardcoded allowlists and one resource-cap rule, all in
operator-side configuration. A user cannot read the lists, cannot appeal an exclusion, and cannot
determine from the chain why one application renews for free and another does not; no appeal path
exists in code \srccode{appsmonitor}{targeted search, no match}.

\item \textbf{On payment failure the only remedy is email.} Construction~\ref{alg:prepay}, failure
mode (a): a message to the application's contact address, rate-limited to one per \qty{24}{\hour} per
subscription, with no in-band signal, no user override, and no path that revives an expired
application \srccode{appsmonitor/\allowbreak{}src/\allowbreak{}services/\allowbreak{}appsSubscriptionRenewalMonitor.js}{19,463--575}.

\item \textbf{The record of what was paid does not survive a week.}
$T_{\mathrm{rec}} = \qty{6}{\day}$ on every collection \srccode{appsmonitor/\allowbreak{}config/\allowbreak{}default.js}{33}.
The on-chain leg is permanent and public; the association between it and a card charge is neither.

\item \textbf{Every owner-keyed capability runs through the custodian.} Because
$\mathrm{addr}(pk^{\mathrm{id}}_u)$ is the specification's \texttt{owner}
(Definition~\ref{def:facilitator}), update, expire, and any future mechanism that reads the owner
field are exercisable by $\mathcal{F}$ and not by $u$.

\item \textbf{A lapse destroys the data, and only the operator can avert it.} Application data at
expiry is always deleted: expiry removes the application from the network and nothing --- volumes
included --- is retained past $e_\app$ \srcintent{design intake}. That behaviour is
the network's and is the same on Part~A, where a tenant can always avert it by signing one more
message and paying for it. Here they cannot: a renewal requires a signature under
$pk^{\mathrm{id}}_u$ (Construction~\ref{alg:settle}, step~2), which only $\mathcal{F}$ can produce
(Property~\ref{prop:custody}), so a declined card and an unread dunning email are together sufficient
to destroy a tenant's data with no action the tenant can take. That composition of item~3 with
deletion, rather than deletion itself, is the trust surface.
\end{enumerate}

\noindent The deletion behaviour in item~6 is stated by the team as current and unconditional
\srcintent{design intake}; no file in this extraction was read that establishes it
independently, so it carries intent provenance and is not tagged \shipped. It is nonetheless a
definite answer and it replaces the \emph{Unclear} this list previously carried. Note that
Invariant~\ref{inv:datasurvives} (``data survives suspension'') remains a property of the
unimplemented layer of Section~\ref{sec:paying:creditspec} alone and must not be read across: on
every deployed path, expiry deletes.

\paragraph{A consequence for Section~23.}
This is the one place where a service outside the protocol leaves a mark inside it. Because
$\mathrm{addr}(pk^{\mathrm{id}}_u)$ is written into the specification's \texttt{owner} field
(Definition~\ref{def:facilitator}), some fraction of the network's nominal application owners are
operator-held keys standing for many unrelated users rather than end users. The measured owner
distribution is consistent with this: one owner identity holds \num{246} of the \num{1195} deployed
applications, \qty{20.6}{\percent} of the set
\srcmeasured{api.runonflux.io/\allowbreak{}apps/\allowbreak{}globalappsspecifications}{2026-09-03}. The measurement cannot
distinguish a custodian from a single large tenant, and this paper does not claim which it is. What
follows regardless is a constraint on every mechanism keyed on ``the owner'': the
owner-concentration statistics of Section~23, moderation standing, and the delegation object of
Section~17 all read a field that may name a service rather than a party, and each must say which it
means. For Section~17 the resolution is already fixed by Corollary~\ref{cor:nodelegation}: wherever
that field names an operator, there is no principal-held key, and the delegation object does not
apply rather than applying weakly.

\paragraph{Whose service this is.}
The framing used throughout this subsection is the team's own, and it is both cleaner and more
accurate than treating custody as a \Flux option with a caveat: \textbf{the user never sees the key;
it is server-held in trust; and this is a centralized service run by a third-party operator, not a
feature of the chain} \srcintent{design intake}. The paper adopts it without apology
and without criticism, because it is what makes the central claim precise rather than defensive.
Part~A is the protocol's settlement path: trustless, shipping today in two independent forms, and an
autonomous agent uses it and meets none of the trust surface above --- no account, no card, no email
address, no operator (Property~\ref{prop:nostanding}, Invariant~\ref{inv:pathindep}). What the
service buys is card payment and the removal of key management for humans; it buys no access that
direct settlement does not already give, and by Property~\ref{prop:chainindifferent} the network
cannot tell whether it exists. The team describes it as a good middle step until an on-chain
mechanism exists \srcintent{design intake}.

The observable direction of travel inside the operator's own stack is currently toward consolidating
custody rather than away from it, and the paper records that beside the stated intent rather than in
place of it: one identity per account, one unified read-only balance view, and a key-provider
abstraction that hardens how the held key is protected without removing custody \inprogress
\srccode{console-api/\allowbreak{}src/\allowbreak{}lib/\allowbreak{}crypto/\allowbreak{}keys.ts}{63--286}
\srccode{console-api/\allowbreak{}src/\allowbreak{}routes/\allowbreak{}billing.ts}{63--165}.

\begin{remark}[The operator's option, which is not a protocol matter]
\label{rem:operatoroption}
The operator may choose to expose $K_u$ to users --- that is, to stop being a custodian --- and
states that it may \srcintent{design intake}. That would be a change to one
company's product and to nothing else: no consensus rule, no \FluxOS release and no \appspec field
would move, and by Property~\ref{prop:chainindifferent} the network could not observe the difference
either before or after. This paper records the option, attributes the decision to the operator, and
does not track it. Nothing in Part~A, in Section~17 or in Section~19 depends on which way it goes.
\end{remark}

\settlednote{A non-custodial construction offering the same convenience is not designed here, and
after \Cref{rem:no-protocol-renewal} it is not a gap in this paper's scope: the protocol offers no
recurring authorisation at all, so there is no protocol object a non-custodial alternative would be an
alternative \emph{to}. What is not open, and is worth restating: the direct settlement path of
Part~A is already non-custodial, already deployed, and already usable by an autonomous agent. What
custody buys is card payment and the removal of key management for humans --- convenience, not
capability (Property~\ref{prop:chainindifferent}).}

\subsubsection{The credits layer, and why it is a service object rather than a protocol one}
\label{sec:paying:creditspec}

\textbf{Nothing in this subsection is a consensus mechanism.} By \Cref{rem:balance-location} the
deposited balance lives in a third-party operator's backend, and by \Cref{rem:no-protocol-renewal}
the protocol offers no automatic renewal to attach it to. What follows is therefore the
specification an \emph{operator} would implement --- the state it carries, the transitions it takes,
and the cost of the one guarantee that distinguishes a credit balance from plain expiry --- and not
a proposal for what nodes should evaluate. It is retained at this depth for two reasons: the
lifecycle is genuinely intricate and no public description of it exists, and a reader evaluating such
a service should be able to check it against a written specification. Every statement is \planned and
written in the conditional, and none of it describes the facilitator of
Section~\ref{sec:paying:facilitator}, which holds no balance and has no lifecycle state
(Property~\ref{prop:nobalance}).

\begin{definition}[Credit state --- \planned, no implementation exists]
\label{def:creditstate}
For a tenant $u$, let $b_u \in \mathbb{Z}_{\ge 0}$ denote a deposited balance in $\sat$; no deployed
system maintains such a quantity (Property~\ref{prop:nobalance}). For each
application $\app$ owned by $u$, let $e_\app \in \mathbb{N}$ denote the paid-through height as in
Definition~\ref{def:expiry} and let
\[
\mathrm{st}(\app, \hgt) \;\in\; \{\textsf{funded},\ \textsf{grace},\ \textsf{suspended},\
\textsf{evicted}\}
\]
denote the lifecycle state at height $\hgt$. Let $\price^{\mathrm{ren}}(\app,\hgt) \defeq
\price(\app, \ninst_\app, \dur_\app, \hgt)$ be the cost of one renewal period under
Definition~\ref{def:price}, $\price^{\mathrm{res}}$ a resumption charge, and $G_{\mathrm{g}},
G_{\mathrm{s}}, \Delta_{\mathrm{pre}} \in \mathbb{N}$ the grace window, suspension window and
pre-authorisation lead, all in blocks. Transitions would be as in Table~\ref{tab:transitions}; every
transition would be evaluated by the operator from chain state plus its own balance store
(\Cref{rem:balance-location}).
\end{definition}

\begin{table}[H]
\centering\small
\caption{Specified transitions of $\mathrm{st}(\app,\hgt)$. \planned --- no implementation exists.}
\label{tab:transitions}
\begin{tabular}{@{}llp{0.27\linewidth}p{0.19\linewidth}p{0.15\linewidth}@{}}
\toprule
from & to & condition & effect on workload & effect on data \\
\midrule
\textsf{funded} & \textsf{funded} &
$\hgt = e_\app - \Delta_{\mathrm{pre}}$ and $b_u \ge \price^{\mathrm{ren}}$; then
$b_u \mathrel{-}= \price^{\mathrm{ren}}$, $e_\app \mathrel{+}= \dur_\app$ &
none & none \\
\textsf{funded} & \textsf{grace} &
$\hgt \ge e_\app$ and $b_u < \price^{\mathrm{ren}}$ &
keeps running & retained \\
\textsf{grace} & \textsf{funded} &
$b_u \ge \price^{\mathrm{ren}}$ while $\hgt \le e_\app + G_{\mathrm{g}}$ &
none & none \\
\textsf{grace} & \textsf{suspended} &
$\hgt > e_\app + G_{\mathrm{g}}$ &
containers stopped, not removed & \textbf{retained}, volumes intact \\
\textsf{suspended} & \textsf{funded} &
$b_u \ge \price^{\mathrm{ren}} + \price^{\mathrm{res}}$ while
$\hgt \le e_\app + G_{\mathrm{g}} + G_{\mathrm{s}}$ &
restarted in place & intact \\
\textsf{suspended} & \textsf{evicted} &
$\hgt > e_\app + G_{\mathrm{g}} + G_{\mathrm{s}}$ &
removed & \textbf{destroyed} \\
\bottomrule
\end{tabular}
\end{table}

The design commitment that matters here is the separation of \emph{stopping} from \emph{deleting}.
Suspension would halt compute without destroying state, so that a lapse recoverable in hours would
not cost the tenant their data; only eviction, after $G_{\mathrm{g}} + G_{\mathrm{s}}$ blocks, would
destroy anything. That guarantee has a price, and the price is quantifiable in the same units as
everything else in this section.

\begin{property}[Bounded suspension subsidy]
\label{prop:subsidy}
\planned. A suspended application would consume disk but neither CPU nor bandwidth, because its
containers would be stopped rather than removed and its volumes retained
(Table~\ref{tab:transitions}). The revenue forgone by the network over the suspension window
(after grace, during which the workload still runs at full reservation),
valued at the list rates of Definition~\ref{def:price}, would therefore be bounded above by
\[
\Delta^{\mathrm{sub}} \;\le\; \ninst_\app \cdot p_{\mathrm{hdd}}(\hgt) \cdot
\resvec^{\mathrm{hdd}}(\app) \cdot \frac{G_{\mathrm{s}}}{L^{\ast}(\hgt)}
\quad \text{FLUX},
\]
which is linear in the declared disk footprint and in the window, and independent of the tenant's CPU
and RAM reservation.
\end{property}

\begin{proof}[Proof sketch]
By Table~\ref{tab:transitions}, from the entry into \textsf{suspended} the workload contributes no CPU or
network load once stopped, and the only resource held is the volume, whose declared size is
$\resvec^{\mathrm{hdd}}(\app)$ per instance. Definition~\ref{def:price} prices that footprint at
$p_{\mathrm{hdd}}\resvec^{\mathrm{hdd}}$ per default lifetime $L^{\ast}$, so a window of
$G_{\mathrm{s}}$ blocks costs that fraction of it. The bound is an over-estimate
because compute is not consumed during the window, because the surcharge terms of $B$ do not
apply, and because $B$ is quoted per three instances while the bound charges it per instance. It is
exact only in the sense that the network cannot re-let the disk while it is held.
\end{proof}

The operational reading is that the guarantee ``a lapse does not delete your data'' costs the network
a disk reservation for a bounded, computable number of blocks after grace, and nothing else beyond the
full reservation already conceded during grace. That is why the
choice of $G_{\mathrm{s}}$ is a genuine parameter with a trade-off rather than a matter of taste
(Section~\ref{sec:paying:open}).

\subsubsection{What would authorise a recurring charge}

A recurring charge needs an authority that outlives the tenant's presence. Three constructions are
available and their costs differ sharply.

\begin{enumerate}\setlength{\itemsep}{2pt}
\item \textbf{Custody.} The simplest construction, and the one the third-party service uses today ---
in a form that does not even require a balance. It is not a protocol mechanism and the network takes
no part in it (Property~\ref{prop:chainindifferent}): the operator holds the tenant's identity and
payment keys and fronts the on-chain payment (Section~\ref{sec:paying:facilitator}), so no
authorisation object is needed at all, there being nothing to authorise when the authoriser already
holds the key. \texttt{usersToExtend}
separately lets one hardcoded address submit expire-only extensions on behalf of any owner
\srccode{flux/\allowbreak{}ZelBack/\allowbreak{}config/\allowbreak{}default.js}{229}. The cost is the trust surface enumerated in
Section~\ref{sec:paying:facilitator} and a key set whose compromise is a fleet-wide event.
\item \textbf{On-chain allowance.} A consensus-level primitive under which the tenant authorises a
bounded draw. Not available on the \Flux L1, for the reason given in Section~17; it would require an
L1 consensus change, which places it in the same class as the header change of Section~22.
\item \textbf{Standing signature or session key.} The tenant pre-signs a bounded renewal authority:
exactly the delegation object specified in Section~17, with $\agent$ instantiated as a renewal
service and $\scopeset$ restricted to \emph{renew this application at no more than this price}, under
budget $\budget$.
\end{enumerate}

\noindent\textbf{Recommendation: option (3)}, because it makes credits a special case of the
delegation of Section~17 rather than a second, parallel mechanism, and because it is the only option
that removes the custodian rather than relocating it. It inherits the enforcement limits of that
section unchanged --- in particular, a delegation is only as bounded as the party that checks the
bound, and Section~17 states where that check would live and where it currently does not exist at
all. One consequence is worth naming: because payer and owner are already independent in the
deployed protocol (Property~\ref{prop:payersep}), option~(3) plus a third-party payer would reproduce
the card-payment convenience of Section~\ref{sec:paying:facilitator} without reproducing its custody.
The obstacle is the authorisation object, not the payment rail.
\begin{remark}[Recurring authorisation is out of protocol scope --- settled]
\label{rem:no-protocol-renewal}
This paper previously carried the choice among custody, an on-chain allowance and a standing
signature as an open construction blocking a credits specification. It is settled, and settled by
exclusion rather than by selection: \textbf{the protocol offers no automatic renewal at all}. A
tenant renews by signing and paying for each period, or delegates that to a third-party operator who
does it on their behalf \srcintent{design intake} (\planned).

The consequence is worth stating plainly rather than treating as a gap. There is no consensus-level
allowance object, no session-key construction in the protocol, and no standing authority a node
evaluates. The recurring-payment problem is moved wholly into the service layer described in
\S\ref{sec:paying:facilitator}, where it is solved by custody --- and where its trust surface is
already enumerated. This is a smaller protocol, not a deferred one, and it removes an entire class of
consensus change from the target architecture. What it costs is stated in
\Cref{rem:agent-must-hold-keys}: an autonomous tenant that wants to renew without a human must hold
a key and sign, or accept a custodian.
\end{remark}

\subsubsection{Metering, billing accuracy and dispute}

The central design fact is easy to state and easy to underrate: \textbf{\Flux prices reservation, not
consumption.} The charge of Definition~\ref{def:price} is a function of the \emph{requested} resource
vector and duration, never of what the workload actually used. There is consequently no metering on
the application path at all --- no counter to read, no counter to attest, and no counter to dispute.
This is the same reasoning that Section~7 applies one layer down: PNR pays all nodes through the PoN
rotation rather than the hosting node precisely so that settlement carries no metering or attestation
dependency (requirement P7 there). The two choices are the same choice made twice, and the
consequence is worth stating plainly: reservation pricing eliminates an entire class of
operator-versus-tenant dispute, because there is no usage record for the two parties to disagree
about.

What remains disputable is a single binary fact --- \emph{did the application run} --- and that fact
is observable by third parties: \FluxOS location records, \FDM health checks, and the tenant's own
probes all speak to it independently. The dispute surface is small by construction. This paper does
not invent an arbitration mechanism, because none exists and none is required to make the above true.

Where metering does exist in the wider stack, it is already inexact, and any future
consumption-priced product would inherit the problem.

\begin{property}[FluxAI metering overdraft]
\label{prop:overdraft}
\shipped. FluxAI gates each request on $\texttt{api\_credit} > 0$ and returns HTTP 402 otherwise, but
deducts usage asynchronously: increments accumulate in Redis under a per-user key with a TTL of
\qty{300}{\second}, and a separate scheduled command applies them to the durable balance inside a
locked transaction and deletes the key
\srccode{fluxai/\allowbreak{}api/\allowbreak{}v1/\allowbreak{}api\_client.py}{1210--1214}
\srccode{fluxai/\allowbreak{}api/\allowbreak{}cache.py}{9--32}
\srccode{fluxai/\allowbreak{}api/\allowbreak{}management/\allowbreak{}commands/\allowbreak{}process\_token\_usage.py}{41--99}.
Two consequences follow. (i) A caller can consume beyond their balance for the whole interval between
flushes, since the gate reads only the durable balance --- a real overdraft window. (ii) If the flush
interval ever exceeds the key TTL, the accumulated usage expires and is never billed at all. The
bound on (i) is the flush cadence, which is \emph{not} determinable from source: the dispatcher runs
per minute, but the job's schedule is stored in the database rather than in the repository
\srcdoc{fluxai/setup/crons/django\_cron:1}.
\end{property}

\begin{proof}[Proof sketch]
The admission check and the deduction read and write different stores, and the only synchronisation
between them is the scheduled flush; therefore the interval between flushes bounds the divergence.
The TTL is set on the accumulator key, not on the durable row, so expiry of the key discards the
pending debit rather than deferring it. Both follow directly from the two cited code paths; neither
requires an adversarial caller.
\end{proof}

The honest statement for a future consumption-priced product is therefore that the overdraft bound is
a scheduling parameter, not a cryptographic one, and that it must be pinned in code before it can be
quoted.

\subsubsection{Service levels and refunds}

No refund or credit mechanism exists in any repository in this corpus. The deployed behaviour on node
failure is \emph{redeployment}, not compensation: when an instance disappears, other nodes observe
that the application is below its declared instance count and one of them installs it, subject to the
broadcast-and-wait collision rule
\srccode{flux/\allowbreak{}ZelBack/\allowbreak{}src/\allowbreak{}services/\allowbreak{}appLifecycle/\allowbreak{}appSpawner.js}{152--247,666--703}; a node that loses
confirmation has its locations evicted from global state by every other node independently
\srccode{flux/\allowbreak{}ZelBack/\allowbreak{}src/\allowbreak{}services/\allowbreak{}appMonitoring/\allowbreak{}nodeStatusMonitor.js}{78--141}. Stated plainly:
\textbf{the remedy is replacement, not refund.}

This is a coherent position rather than an omission, and the reason is the one above. An SLA credit
is a function of measured availability, so paying one would require exactly the metering the design
has deliberately avoided --- and it would reintroduce the attestation dependency that Section~7
removed. Promising a refund path that nothing implements would be worse than saying so.
\begin{remark}[No protocol availability remedy --- settled]
\label{rem:no-sla}
Whether the target architecture should acquire an availability remedy --- a refund or credit when a
node fails to serve a paid deployment --- is settled: \textbf{it should not, at the protocol layer}.
Any such remedy is a service-layer product, offered commercially by an operator if at all
\srcintent{design intake} (\planned).

This is the consistent choice rather than a gap, and it follows from a fact this section already
establishes: \Flux{} prices \emph{reservation}, not consumption (\S\ref{sec:paying:open}). A
tenant buys reserved capacity, so there is no delivered-uptime quantity for a remedy to be computed
against, and no third-party-observable availability signal exists from which one could be. Building
that signal would mean building the attestation network first --- which \S\ref{sec:feasibility}
classes as Research --- and then defining an oracle that decides when a deployment was unserved.
Deciding not to have a protocol remedy avoids inventing both.

It composes with the settled positions on credits and renewal (\Cref{rem:balance-location},
\Cref{rem:no-protocol-renewal}): convenience, custody and commercial guarantees all sit in the
service layer, and the protocol keeps one job --- deciding whether a deployment was paid for.
\end{remark}

\subsubsection{Invariants}
\label{sec:paying:invariants}

Stated for all three objects --- the protocol's direct settlement (Part~A), the operator's deployed
service (Section~\ref{sec:paying:facilitator}), and the unimplemented credits layer
(Section~\ref{sec:paying:creditspec}). Where an invariant holds on one and not another, that is said.
Only the first is a claim about \Flux; the second is a claim about one company's product and binds
nobody else who builds over Part~A.

\begin{enumerate}\setlength{\itemsep}{3pt}
\item\label{inv:conservation} \textbf{Conservation.} $b_u$ decreases only by renewals actually
applied, and every debit corresponds to exactly one extension of exactly one $e_\app$ and, after
$\hpa$, to exactly one PNR transaction on L1 (Requirement~\ref{req:pnr}). No debit exists without an
on-chain settlement; no settlement exists without a debit. \planned as stated, since $b_u$ has no
deployed referent (Property~\ref{prop:nobalance}). On Part~A it holds trivially, the payment and the
extension being the same event; on the operator's service it holds in the degenerate form that each card
charge is settled by exactly one transaction to the sink and no ledger stands between them
(Construction~\ref{alg:prepay}, step~4).

\item\label{inv:nosilent} \textbf{No silent expiry.} Every transition out of \textsf{funded} shall be
observable by the tenant before it takes effect. On Part~A the \emph{expiry} half holds by
Property~\ref{prop:nostanding} --- $e_\app$ is public and computable in advance by anyone --- but the
\emph{entry} half fails: an underpaid registration never becomes \textsf{funded} and the tenant is
not told (Property~\ref{prop:silentunderpay}). Requirement~\ref{req:nack} is what would repair it,
and it is a requirement on every path. On the operator's service the impending lapse \emph{is} signalled,
but by email only, on the operator's schedule, at most once per \qty{24}{\hour}, with no in-band
channel and no user override (Construction~\ref{alg:prepay}, failure mode (a)); that is weaker than
this invariant asks and is recorded as such rather than counted as satisfying it.

\item\label{inv:datasurvives} \textbf{Data survives suspension.} \planned, and a property of
Section~\ref{sec:paying:creditspec} alone. No transition into \textsf{suspended} destroys tenant
data; only \textsf{evicted} does, and only after $G_{\mathrm{g}} + G_{\mathrm{s}}$ blocks. By
construction of Table~\ref{tab:transitions}; the cost is bounded by Property~\ref{prop:subsidy}. The
deployed network has no suspension state at all, and application data at expiry is always deleted
\srcintent{design intake}; this invariant must not be read across to it. The
contrast is the point: on every deployed path expiry is destructive, so ``data survives suspension''
would be a guarantee a credits layer \emph{introduces}, not one it preserves.

\item\label{inv:pathindep} \textbf{Path independence.} For every deployment reachable through the
operator's service or through a credits layer, the same deployment is reachable through direct
payment with an explicit transaction.
\emph{Proof.} \textbf{Deployed case.} By the invariant of Construction~\ref{alg:prepay}, the
facilitator's whole action reduces to (a) a signature over $m_\app$ under the key named in the
specification's \texttt{owner} field and (b) one transaction paying $\price$ to the sink carrying
$H(m_\app)$ --- that is, exactly steps~2 and~3 of Construction~\ref{alg:settle}. A tenant holding
their own key performs both and obtains the identical $e_\app$ by Definition~\ref{def:expiry}.
\textbf{Specified case.} A credit-path renewal at height $\hgt$ debits
$\price^{\mathrm{ren}}(\app,\hgt)$ and extends $e_\app$ by $\dur_\app$; by
Requirement~\ref{req:pnr} it does so through one L1 settlement of that amount, which the tenant may
broadcast themselves at the same height for the identical $e_\app$, since the promotion test depends
only on $(v^{\mathrm{obs}}, \hgt, \app)$ and on no property of the payer
(Property~\ref{prop:payersep}). Hence the reachable sets coincide in both cases. $\square$
\textbf{Neither the operator's service nor a credits layer adds capability}, and a user who uses
neither loses none. Together with Property~\ref{prop:chainindifferent} --- the network cannot see the
service --- this is what places the service strictly on top of \Flux rather than inside it: it adds
no reach downward and imposes no requirement upward. It is the load-bearing item in this list, and
the reason the framing of Section~\ref{sec:paying:facilitator} is a description rather than a
defence.

\item\label{inv:monotone} \textbf{Monotone expiry.} $e_\app$ is non-decreasing while
$\mathrm{st}(\app,\hgt) \ne \textsf{evicted}$. Immediate from Table~\ref{tab:transitions}, in which
the only writes to $e_\app$ are increments. Note that the deployed system does \emph{not} satisfy the
analogous property under the expire-only update path, where a Foundation-held key can submit a
\emph{reducing} expire value with no ``new expire must be larger'' check
\srccode{appsmonitor/\allowbreak{}fixOverExtendedAppExpiry.js}{1--24}; that is an operational tool, and any
credits implementation must not inherit the exemption.
\end{enumerate}

\subsubsection{Open parameters}
\label{sec:paying:open}

These are service-layer parameters by \Cref{rem:balance-location}: an operator chooses them, and
nothing in consensus depends on them. The paper nonetheless recommends values, because an operator
building this needs a defensible starting point and because the trade-offs are the same whoever
implements them. Each was proposed by this paper and has been confirmed by the author
(Table~\ref{tab:credit-params}).

\begin{table}[H]
\centering\footnotesize
\caption{Lifecycle parameters for a credits operator. Proposed by this paper and \textbf{confirmed
by the author} \srcintent{design intake}. Block counts assume the
\SI{30}{\second} post-\PoN target, i.e.\ \num{2880} blocks per day
\srccode{fluxd/src/chainparams.cpp}{105}. \planned.}
\label{tab:credit-params}
\begin{tabular}{@{}llll@{}}
\toprule
Parameter & Domain & \textbf{Recommended} & Reasoning \\
\midrule
$G_{\mathrm{g}}$ grace & $[240, 2880]$ & \textbf{\num{2880}} (\SI{24}{\hour}) & covers a full
business day \\
$G_{\mathrm{s}}$ suspension & $[2880, 43200]$ & \textbf{\num{20160}} (\SI{7}{\day}) & one week of
stored, unpaid data \\
$\Delta_{\mathrm{pre}}$ pre-auth lead & $[20, \infty)$ & \textbf{\num{120}} (\SI{1}{\hour}) &
$\gg$ confirmation depth \\
$\price^{\mathrm{res}}$ resumption & $[0, \infty)$ & \textbf{one renewal period} & prices the
subsidy exactly \\
\bottomrule
\end{tabular}
\end{table}

\paragraph{Why these values.}
$G_{\mathrm{g}} = \num{2880}$ blocks, one day, is chosen so that a funding failure occurring at any
hour is recoverable within one business day without the tenant losing service. The shorter end of the
domain (\SI{2}{\hour}) converts an overnight bridge delay into a suspension, which is the routine
failure this window exists to absorb; the cost of the longer setting is one day of unpaid load at
full occupancy, i.e.\ $G_{\mathrm{g}}/L^{\ast}$ of one renewal period, which lies outside the disk-only
bound of Property~\ref{prop:subsidy}.

$G_{\mathrm{s}} = \num{20160}$ blocks, seven days, is the willingness-to-subsidise statement. The
question this parameter really asks is how long the network stores data nobody is paying for. A week
is long enough that a tenant returning from an outage, a holiday or a payment-method failure finds
their deployment intact, and short enough that the storage subsidy stays bounded at one week per
evicted application rather than accumulating indefinitely. Note that the fifteen-day end of the
domain doubles the subsidy to buy little: recovery that has not happened within a week is
overwhelmingly abandonment rather than delay.

$\Delta_{\mathrm{pre}} = \num{120}$ blocks, one hour, sits an order of magnitude above the binding
constraint. The floor is confirmation depth plus specification propagation; propagation in
Construction~\ref{alg:settle} is seconds and confirmation dominates. One hour leaves room for a
congested mempool and a re-sent transaction without approaching the upper risk, which is that a
renewal authorised too far ahead is priced against a stale segment of
Definition~\ref{def:price}.

$\price^{\mathrm{res}}$ equal to \emph{one renewal period at the application's own price} is the
only anchor that is correct by construction rather than by taste. Property~\ref{prop:subsidy} prices
the suspension window as a subsidy; charging exactly one period on resumption returns that subsidy
from the tenant who consumed it, scales automatically with the resources the application actually
reserved, and needs no separate schedule to maintain. Zero would make the suspension window free cold
storage, which is the abuse the property identifies; a flat fee would misprice small and large
applications in opposite directions.

\begin{remark}[Where a deposited balance lives --- settled: nowhere on the chain]
\label{rem:balance-location}
The location of a deposited balance was the architectural choice this section could not make. It is
now settled, and the answer is that the object is not a protocol object at all. Credit balances are
held \textbf{off-chain, per user, in the backend of a third-party operator} --- the team names InFlux
Technologies, building on Flux Cloud, as such an operator. That service is account-based: a user may
sign in with SSO or email, the backend derives a ZelID for them, they are logged in under it, and the
operator holds a per-user balance which the user tops up and authorises spend from
\srcintent{design intake} (\planned).

Three consequences follow, and the paper states them rather than the convenience.
\begin{enumerate}
  \item \textbf{Invariant~\ref{inv:conservation} is enforced by the operator's bookkeeping.} Not by
  consensus, not by a node quorum, not by the tenant's device. A credit balance is a claim against an
  operator, redeemable at that operator's discretion and subject to its solvency.
  \item \textbf{The chain cannot see it, and by Property~\ref{prop:chainindifferent} does not need
  to.} What reaches the chain is an ordinary payment for a deployment, indistinguishable from one a
  tenant signs directly. Credits are therefore invisible to \PNR, to the fee pool, and to every
  property proved in \S\ref{sec:pnr}.
  \item \textbf{The specification below is retained as a description of the service-layer object},
  not as a protocol proposal. Its state machine is what such an operator must implement correctly;
  it is not what nodes evaluate.
\end{enumerate}
\end{remark}

\begin{remark}[\specv{9} replaces \texttt{blocksLasting} with \texttt{duration} --- settled from source]
\label{rem:v9-duration}
This was carried as an open question. It is answered by the schema itself: in the \specv{9} schema
\texttt{duration} is a \textbf{top-level required} property --- \texttt{\{"type":"integer",\allowbreak{}
"minimum":1\}}, documented as ``requested deployment duration in seconds (first-class in v9; the
foundation the subscription and auto-renew work builds on)'' --- and \textbf{\texttt{blocksLasting}
does not appear in the \specv{9} schema at all}, nor does \texttt{expire}
\srcbranch{flux}{feat/appspec-v9}{ZelBack/src/services/appRequirements/schemas/v9.json}.

Two consequences follow. \emph{First}, $L^{\ast}(\hgt)$ in Definition~\ref{def:price} becomes a unit
conversion rather than a fork-dependent constant, and the post-\PoN{} $4\times$ recomputation of
Definition~\ref{def:expiry} becomes dead code on the \specv{9} path --- lifetime is expressed in
seconds, so it no longer moves when block spacing changes. That is a simplification, and it removes
a class of error the $4\times$ branch has already caused once (\S\ref{sec:pnr}).

\emph{Second, and this is the caveat that keeps the claim honest}: the field is declared and required
but \textbf{not yet consumed}. A search of \texttt{app\allowbreak{}Requirements/\allowbreak{}} and \texttt{appLifecycle/} on the
same branch finds \texttt{duration} only in comments; \texttt{blocksLasting} is still what the
deployed v8 install path reads \srcbranch{flux}{feat/appspec-v9}{ZelBack/src/services/appLifecycle/appInstaller.js}.
So \specv{9} \emph{specifies} the replacement and does not yet \emph{implement} it --- the same
declared-but-inert pattern \S\ref{sec:one-platform} documents for \texttt{paymentMemo},
\texttt{referralCode} and \texttt{machineType}, and the reason \Cref{rem:v9-timing} constrains
$\hstar$ to follow the implementation rather than precede it.
\end{remark}

\subsubsection{Comparison}

\textbf{Akash} escrows a lease account and streams payment per block, closing the lease when the
escrow empties \citep{akash}; that is the closest analogue to the credits layer specified in
Section~\ref{sec:paying:creditspec}, with the balance on-chain and the state machine in consensus
rather than in an orchestrator or in a custodian's database. Its answer to ``where does the balance
live'' is the opposite of the one settled here --- off-chain, in an operator's backend
(\Cref{rem:balance-location}) --- and it pays for that answer with a consensus-level
account primitive the \Flux L1 does not have. \textbf{Filecoin} deals prepay and collateralise, with
provider slashing for failure \citep{filecoin2017}; \Flux has no equivalent slashing on the
application path, which is precisely why its remedy is replacement rather than compensation --- there
is no bonded party to take value from. \textbf{Hyperscaler billing} is post-paid,
consumption-metered, and correspondingly dispute-heavy; reservation pricing is the deliberate
opposite, trading utilisation efficiency for verifiability, and Property~\ref{prop:overdraft} is a
small demonstration of what the metered path costs even inside \Flux's own stack.
\textbf{Lightning invoices} are the closest match to Part~A's ``the payment is the entitlement''
property \citep{poon2016lightning}: a payment preimage is the proof of settlement and no account
mediates it. \CumulusVPN's memo-derived entitlement is effectively an on-chain invoice with a public
settlement record --- which is exactly the property that makes it work without a server and exactly
the property that makes it the problem statement for Section~21.

%% Drain this section's float queue before the next begins.
\FloatBarrier
